\section{DPO vs PPO：实证探索}

\subsection{实验设计}

Lambert 的团队在 AI2 进行了一系列系统实验，试图回答：\textbf{PPO 是否真的比 DPO 更好？如果是，好在哪里？}

实验基于 Llama 2 13B 模型，逐步叠加不同的训练方案：

\begin{figure}[H]
\centering
\includegraphics[width=0.95\textwidth]{slides-images/slide-062.jpg}
\caption{起点：Llama 2 Base（蓝色）经指令微调后（红色）的性能提升\protect\footnotemark}
\end{figure}
\footnotetext{Slides 第63页。}

\begin{importantbox}{指令微调的提升幅度远大于偏好优化}
在所有实验中，从 base 模型到 SFT 模型的\textbf{指令微调}阶段带来的提升最大——它将模型"放到了地图上"。之后的 DPO/PPO 阶段带来的提升通常在 0--2\% 之间。这提醒我们：对齐优化的收益是边际的，但正是这些边际改进决定了模型在竞争中的位置。
\end{importantbox}

\subsection{数据集的影响}

\begin{figure}[H]
\centering
\includegraphics[width=0.95\textwidth]{slides-images/slide-065.jpg}
\caption{不同数据集的 DPO 训练效果对比\protect\footnotemark}
\end{figure}
\footnotetext{Slides 第66页。}

实验中测试了多个偏好数据集：

\begin{itemize}
    \item \textbf{Anthropic HH-RLHF}：虽然噪声较大，但仍能提供小幅提升
    \item \textbf{UltraFeedback}：提供了更大的性能提升
    \item 其他多个开源数据集：效果参差不齐
\end{itemize}

一个令人意外的发现是：\textbf{在事实性（factuality）维度上，几乎所有偏好数据集都不起作用}。这说明当前的偏好优化方法在提升模型知识准确性方面存在根本限制。

\begin{figure}[H]
\centering
\includegraphics[width=0.95\textwidth]{slides-images/slide-072.jpg}
\caption{多个偏好数据集的训练效果对比：UltraFeedback 始终是最优选择\protect\footnotemark}
\end{figure}
\footnotetext{Slides 第73页。}

\subsection{PPO 的实验结果}

当团队实现并测试 PPO 后：

\begin{figure}[H]
\centering
\includegraphics[width=0.95\textwidth]{slides-images/slide-067.jpg}
\caption{从 DPO 切换到 PPO：在多个数据集上 PPO 略优于 DPO（约 1\%）\protect\footnotemark}
\end{figure}
\footnotetext{Slides 第68页。}

\begin{itemize}
    \item PPO 在大多数情况下\textbf{略优于 DPO}，大约 1\% 的提升
    \item 使用更大的奖励模型（70B）并没有带来预期的显著提升
    \item 添加更多类型的提示（代码、推理）对特定评估有帮助，但不改变整体平均分
\end{itemize}

\begin{figure}[H]
\centering
\includegraphics[width=0.95\textwidth]{slides-images/slide-070.jpg}
\caption{PPO 的实验总结\protect\footnotemark}
\end{figure}
\footnotetext{Slides 第71页。}

\begin{importantbox}{PPO 的关键发现}
\begin{enumerate}
    \item \textbf{"总有下一个要消融的东西"}：PPO 的超参数空间极大（正则化、value function、warmup、batch size 等），每次调好一个还有下一个
    \item \textbf{"PPO 能得到最好的模型，但我们不知道为什么"}：PPO 可靠地给出了稍好的结果，但缺乏对其优势来源的清晰理解
    \item \textbf{生成是最大瓶颈}：PPO 需要在训练过程中不断从策略模型生成新回复，这比 DPO 慢得多
\end{enumerate}
\end{importantbox}

\begin{warningbox}{PPO 的投入产出比令人质疑}
Lambert 坦言：PPO 虽然可靠地产出了稍好的模型，但投入的工程和计算资源远超 DPO。对于一个只有几名研究生的学术团队来说，"不知道这是否值得"。他理解为什么 OpenAI 使用 PPO——因为在他们的规模上，1\% 的改进意义重大——但对于大多数研究者来说，DPO 的投入产出比更好。
\end{warningbox}

\subsection{本章小结}

通过系统的实证比较，Lambert 的团队得出了几个重要结论：（1）数据质量比算法选择更重要——UltraFeedback 在 DPO 和 PPO 上都显著优于其他数据集；（2）PPO 略优于 DPO，但代价是巨大的工程复杂性；（3）更大的奖励模型不一定带来更好的最终策略；（4）对齐优化的边际收益递减——从 base 到 SFT 的提升远大于从 SFT 到 DPO/PPO 的提升。

%% ========== Part 6: Online 方法 ==========

